\scanpage{137}
\mspage{8}
% archive sheet number: 136

À supposer qu'une telle construction existe dans un cas tant soit peu général, elle s'annonce un \unsure{tantinet} délicate. Pour fixer les idées, il serait bon \deleted{\unsure{tjs|de}} au moins de regarder déjà les cas suivants

$1^\circ)$ Cas d'un atelier ensembliste

\deleted{\unsure{$2^\circ)$}} Déjà \unsure{ça} n'a pas l'air évident, surtout pas quand on ne suppose pas que l'atelier soit modéré, et $\unsure{\eqsim|\simeq|=}$ pour des $N$ très particuliers comme p.ex. $\operatorname{Omb}^{\unsure{\bullet|\circ}}F$, $F$ une figure, ou $\operatorname{Supp}F$.

$2^\circ)$ Cas d'une maquette $A$, $N$ partie quelconque de $A$ (NB $x\overset{\circ}{\ll}y\Leftrightarrow x=y$ !). Là, on doit trouver la maquette $A\setminus N$ simplement. Ce n'est \unsure{pas} \unsure{sorcier} cette fois !

\par\noindent\hspace*{0.3\linewidth}\rule{0.25\linewidth}{0.4pt}\par

Pour la relation avec le \unsure{formalisme} des supports, s'inspirer éventuellement de \unsure{GF|CF} $\overline{\text{\uline{VIII}}}$ p.~70 ff. Là je n'ai pas la pêche pour, et laisse tomber pour le moment.
